\documentclass[journal,10pt]{IEEEtran}
\usepackage{amsmath,amssymb}
\usepackage{graphicx}
\usepackage{booktabs}
\usepackage{cite}
\usepackage[hidelinks]{hyperref}

\begin{document}

\title{Span-Supervised Transformers for Remaining Useful Life Estimation:\\ Train on Two Spans, Infer on One}

\author{Anonymous Authors\thanks{Manuscript submitted for review.}}

\markboth{IEEE Transactions on Instrumentation and Measurement}{Anonymous \MakeLowercase{\textit{et al.}}: Span-Supervised Transformers for RUL Estimation}

\maketitle

\begin{abstract}
Remaining useful life (RUL) estimators are usually trained on a fixed-length
window, yet engines in the same fleet carry histories of very different lengths:
in the C-MAPSS test sets, between 21\% and 35\% of engines have fewer than 100
recorded cycles. The usual remedy---left-pad to a fixed length and pool over
everything---silently mixes padding with measurement. This paper asks a narrower
question than most architecture papers: \emph{given a shared encoder, which span}
\emph{of the history should be pooled, and should several spans be fused?} We answer
with a controlled study on C-MAPSS FD001--FD004. A two-head model supervises one
head on the most recent 50 valid steps and one on the full valid history; at
inference the long-span head alone is used. Across four fleets and fifteen models
(five folds, three seeds) the ensemble attains 9.64, 10.75, 10.68, and 11.19
cycles of test RMSE. This is 7.6\%, 10.6\%, and 20.7\% below a 2026
decision-conditioned causal-graph baseline on FD001, FD002, and FD004
respectively, and 1.9\% above it on FD003. Fusing the two spans at inference is not
neutral: it costs 2.1\%--5.6\%, and a two-branch pipeline costs 18\%--39\%. Training
on two spans while inferring on the longer one is what pays. We also report,
rather than hide, a systematic over-prediction of 1.6--2.4 cycles concentrated at
60--90 cycles of true life.
\end{abstract}

\begin{IEEEkeywords}
Remaining useful life, C-MAPSS, transformer, variable-length sequences, masked
pooling, predictive maintenance.
\end{IEEEkeywords}

\section{Introduction}

Predicting how long a machine will keep running is a regression problem with an
asymmetric cost: over-estimating life delays maintenance, under-estimating it
wastes hardware. Deep models now reach test RMSE near ten cycles on the C-MAPSS
benchmarks \cite{saxena2008}, and the recent literature has moved toward
increasingly elaborate machinery---graph convolutions, causal discovery, state
space models \cite{ast2026,mfssc2024,dkdgjln2026}.

Much less attention has been paid to a humbler design choice. A model sees a
window of the last $L$ cycles. Real histories are shorter than $L$ for a large
minority of engines. The conventional fix is to left-pad and pool over the full
$L$ positions. That works, but it means the pooled representation of a 30-cycle
engine is an average in which 70 of 100 positions are padding. It also raises a
question that architecture papers rarely isolate: if the model could pool over
\emph{all} available history instead of a fixed suffix, would it do better, and
would it help to combine several spans?

This paper isolates exactly that. We hold the encoder, the optimiser, the data
pipeline, and the hyper-parameters fixed, and vary only \emph{which valid positions}
\emph{are pooled} and \emph{whether spans are fused}. The resulting study is narrow but,
we believe, useful: it produces a model competitive with the 2026 state of the art
on three of four fleets, and it produces two negative results that we expect to
save other groups time.

\emph{Contributions.}
\begin{enumerate}
\item A span-supervised formulation in which a single encoder carries two pooling
heads---one over the recent 50 valid steps, one over the full valid history---and
the training loss supervises both while inference reads only the long head.
\item A controlled ablation on FD002 and FD004 showing that the ordering
\texttt{long} $<$ \texttt{dual} $<$ \texttt{short} $<$ \texttt{repeat} is stable
across folds: extra history helps, but \emph{fusing} spans does not.
\item An evaluation on the official C-MAPSS test engines with engine-clustered
bootstrap confidence intervals, reported separately for single models and for the
ensemble, plus an error-versus-life analysis that locates the residual bias.
\end{enumerate}

\section{Related Work}

\textbf{RUL on C-MAPSS.} The benchmark and its scoring function are from
\cite{saxena2008}. Sequence models have largely displaced hand-engineered features;
\cite{mfssc2024} couples multi-scale convolution with a graph over sensors, and
\cite{ast2026} conditions a causal graph on a discrete decision state before a Mamba
backbone. \cite{dkdgjln2026} reports strong FD002 and FD004 numbers with a
domain-generalisation objective. These methods differ from ours in that they change
the encoder; we change only the pooling.

\textbf{Pooling and variable length.} Masked mean pooling over valid positions is
standard in NLP and increasingly used in sensor models. What is not standard is
treating the pooled span as a design variable and measuring it. Work that pads short
histories and pools over the pad implicitly assumes the encoder learns to ignore
padding; our \texttt{long} arm tests that assumption directly.

\textbf{Supervision with auxiliary spans.} Multi-scale supervision is common, but the
usual deployment fuses the scales. Our ablation separates the two decisions:
supervise with several spans (helps) versus fuse them at inference (hurts).

\section{Method}

\subsection{Inputs and window}
Each engine contributes a multivariate sequence of 14 sensor channels and three
operational settings. Windows of $L=100$ cycles are cut with the target
$\min(n-t,{B},125)$, where $n$ is the engine last cycle and $t$ the window end; the
cap follows common practice and reduces the dominance of the long-life tail. A
window that starts before the first cycle is left-padded and carries a binary mask
$\mathbf{m}\in\{0,1\}^{L}$ with $m_i=1$ for valid steps.

\subsection{Encoder}
Let $\mathbf{x}\in\mathbb{R}^{L\times 14}$ and let $\mathbf{b}$ be the first-cycle
baseline used for normalisation. Per-step features are projected to dimension
$d=64$ and passed through a shared temporal convolution with GELU, residual
connection, and layer normalisation. A learned attention over the 14 sensors
produces one token per step,
\begin{equation}
\mathbf{u}_t=\sum_{s=1}^{14}\alpha_{t,s}\,\mathbf{h}_{t,s},
\qquad
\alpha_{t,\cdot}=\mathrm{softmax}(\mathbf{w}^{\top}\mathbf{h}_{t,\cdot}),
\end{equation}
so the sensor mix is allowed to change along the life of the engine.

Positional information is indexed by \emph{distance from the newest step}, not by
absolute position. A learnable table $\mathbf{P}\in\mathbb{R}^{d\times L}$ is
added as $\mathbf{u}_t+\mathbf{P}_{:,L-1-t}$. This makes the representation
translation-equivariant in absolute time: a 40-cycle history and a 100-cycle
history are encoded by the same function of distance-from-now. It is the single
change that lets one encoder serve both spans without a second parameter set.

A transformer encoder of $n_\ell$ layers then mixes the steps. Finally, a
sinusoidal encoding of the absolute cycle index is projected and added, so the
model knows how far into its life the window sits.

\subsection{Two heads, one inference path}
Let $\mathcal{V}=\{i:m_i=1\}$ and let
$\mathcal{S}=\{i\in\mathcal{V}:\,L-1-i<50\}$ be the recent 50 valid steps. With
$\bar{\mathbf{h}}(\mathcal{A})=\frac{1}{|\mathcal{A}|}\sum_{i\in\mathcal{A}}\mathbf{u}_i$,
\begin{equation}
\hat y_A = f_\theta\!\left(\bar{\mathbf{h}}(\mathcal{S})\right),
\qquad
\hat y_B = f_\theta\!\left(\bar{\mathbf{h}}(\mathcal{V})\right),
\end{equation}
where $f_\theta$ is a shared regression head. The training objective is
\begin{equation}
\mathcal{L}
= \ell(\hat y_B, y)
+ \tfrac12\!\left[\ell(\hat y_A,y)+\ell(\hat y_B,y)\right]
+ \tfrac12\,\mathbb{E}\!\left[w(e_B)\,e_B^{2}\right],
\label{eq:loss}
\end{equation}
with $e_B=\hat y_B-y$, $\ell$ the squared error, and $w(e)=1.8$ if $e>0$ and
$1.0$ otherwise. The first term trains the deployed head; the second keeps the
short head informative so that the long head is not the only thing the encoder is
asked to serve; the third prices over-prediction more dearly than
under-prediction, which is the correct asymmetry for maintenance.

At inference we read $\hat y_B$ and discard $\hat y_A$. The ablation in
Section~\ref{sec:abl} is precisely about whether that last sentence is right.

\subsection{Protocol discipline}
Folds are formed over engines, never over windows, so an engine appears in exactly
one fold. Normalisation statistics and the sensor baseline are fitted on the
training folds only. Early stopping uses the held-out fold; the official test set
is read only after a configuration is frozen. We report single-model means and
ensembles in separate rows, as mixing them is the most common way these benchmarks
are inflated.

\section{Experiments}

\subsection{Setup}
C-MAPSS FD001--FD004 \cite{saxena2008}. Five outer folds $\times$ three seeds
(101, 202, 303) give fifteen models per fleet. One transformer layer, $d=64$,
78{,}978 parameters, AdamW, 80 epochs with patience 20 on validation RMSE. No
test-set information enters model selection.

\subsection{Main results}
Table~\ref{tab:main} reports the official-test RMSE. The single-model column is the
mean over fifteen independently trained models; the ensemble column averages their
fifteen predictions per engine.

\begin{table}[t]
\centering
\caption{Official-test RMSE in cycles. Single is the mean $\pm$ standard deviation
over 15 models (5 folds $\times$ 3 seeds); Ensemble averages their predictions. MAE,
NASA score \cite{saxena2008}, and signed bias refer to the ensemble. $\Delta$ is
relative to \cite{ast2026}.}
\label{tab:main}
\small
\begin{tabular}{lrrrrrr}
\toprule
Fleet & Single & Ensemble & MAE & NASA & bias & $\Delta$ \\
\midrule
FD001 & 10.75$\pm$0.67 & \textbf{9.64} & 6.50 & 1.75 & $+1.65$ & $-7.6\%$ \\
FD002 & 11.50$\pm$0.52 & \textbf{10.75} & 6.83 & 2.53 & $+1.96$ & $-10.6\%$ \\
FD003 & 11.69$\pm$1.00 & 10.68 & 6.63 & 2.56 & $+2.40$ & $+1.9\%$ \\
FD004 & 12.12$\pm$0.35 & \textbf{11.19} & 7.12 & 2.61 & $+1.60$ & $-20.7\%$ \\
\bottomrule
\end{tabular}
\end{table}

\begin{figure}[t]
\centering
\includegraphics[width=\columnwidth]{paper_figures_v4/fig_compare.pdf}
\caption{Official-test RMSE per fleet. Grey: the 2026 baseline; blue: our ensemble.}
\end{figure}

Three of four fleets improve on the 2026 baseline, by margins from 7.6\% to 20.7\%.
FD003 is 1.9\% worse and is discussed in Section~\ref{sec:fd003}.

The gap between the single-model mean and the ensemble is 0.6--1.1 cycles. It is
the width of the fold-to-fold distribution, not a tuning artefact: the best and
worst single models on FD003 differ by 3.6 cycles. Any claim on these benchmarks

\begin{figure}[t]
\centering
\includegraphics[width=\columnwidth]{paper_figures_v4/fig_scatter.pdf}
\caption{Ensemble prediction against true RUL on the official test engines. Dashed: y = x.}
\end{figure}
should therefore state which of the two columns it refers to.

\subsection{Ablation: which span, and does fusion help?}
\label{sec:abl}

We hold everything fixed and vary only the pooling and fusion. Four arms:
\texttt{short} pools the recent 50 valid steps; \texttt{long} pools the full valid
history; \texttt{dual} computes a length-dependent convex combination of the two
pooled vectors, $g\,\bar{\mathbf{h}}(\mathcal{V})+(1-g)\,\bar{\mathbf{h}}(\mathcal{S})$
with $g=\mathrm{clip}((|\mathcal{V}|-50)/50,0,1)$, and reads one head from it;
\texttt{repeat} routes the short head output through the long head so that the
two-branch pipeline exists but carries no extra span information.

\begin{table}[t]
\centering
\caption{Ablation. Engine-equal RMSE on held-out folds, mean over two folds.
Percentages are relative to \texttt{long}.}
\label{tab:abl}
\small
\begin{tabular}{lrrrr}
\toprule
Fleet & \texttt{long} & \texttt{dual} & \texttt{short} & \texttt{repeat} \\
\midrule
FD002 & \textbf{9.721} & 10.266 (+5.6\%) & 10.342 (+6.4\%) & 13.493 (+38.8\%) \\
FD004 & \textbf{9.213} & 9.410 (+2.1\%) & 9.678 (+5.1\%) & 10.841 (+17.7\%) \\
\bottomrule
\end{tabular}
\end{table}

Table~\ref{tab:abl} gives an unambiguous ordering,
\texttt{long} $<$ \texttt{dual} $<$ \texttt{short} $<$ \texttt{repeat}, the same
on both fleets and in all four folds. Three readings follow.

\emph{History beyond the last 50 cycles is worth 5--6\%.} That is the difference
between \texttt{short} and \texttt{long}, and it is the direct measurement of what
masked pooling over the full valid span buys.

\emph{Fusing the spans is a net loss of 2--6\%.} \texttt{dual} is strictly worse
than \texttt{long} on both fleets. The gate is not mis-calibrated---it is a
deterministic function of the observable length---so the loss comes from mixing a
dominated estimate into a good one. The short-span head is a useful regulariser
during training and a harmful expert at test time.

\emph{The two-branch pipeline itself is expensive.} \texttt{repeat} carries no
additional span information yet is 18--39\% worse, which is what one expects when a
redundant branch doubles the path length through the network.

\begin{figure}[t]
\centering
\includegraphics[width=\columnwidth]{paper_figures_v4/fig_bias.pdf}
\caption{Mean signed error of the ensemble against true RUL, 25-cycle bins. Positive means over-estimation of remaining life.}
\end{figure}

\subsection{Where the residual error sits}

The ensemble over-predicts by 1.6--2.4 cycles on average, and the bias is not
uniform. It is close to zero for engines near failure and for engines with more
than 100 cycles left, and peaks at $+5$ to $+10$ cycles when the true RUL is
60--90 cycles. That is the interval in which the model has seen enough history to
become confident and not yet enough degradation to be corrected by the data.

This is the honest limitation of the method. The asymmetric penalty in
(\ref{eq:loss}) is meant to counter exactly this, and the NASA score---which
prices over-prediction ten times more heavily---stays between 1.75 and 2.61, so the
tail is controlled even though the mean is optimistic.

\subsection{Why FD003}
\label{sec:fd003}

FD003 is the only fleet where the method does not lead. It is also the fleet with
the largest fold-to-fold spread (standard deviation 1.00 against 0.35--0.67
elsewhere) and the fleet where the two failure modes are drawn from a single
operating condition. Our reading is that the long-span advantage measured in
Table~\ref{tab:abl} is largest where histories are long and operating conditions
vary (FD002, FD004), and that FD003 removes both sources of variation. We report
the number as measured rather than tuning on it.

\subsection{Statistical treatment}

Confidence intervals in Table~\ref{tab:ci} resample \emph{engines}, not engine-fold
pairs, with 5000 bootstrap replicates. Resampling at the engine-fold level
understates the interval by roughly a third because the same engine reappears
across folds.

\begin{table}[t]
\centering
\caption{Ensemble metrics with 95\% engine-clustered bootstrap intervals.}
\label{tab:ci}
\small
\begin{tabular}{lrrr}
\toprule
Fleet & Engines & RMSE & 95\% CI \\
\midrule
FD001 & 100 & 9.636 & [7.575, 11.678] \\
FD002 & 259 & 10.751 & [9.195, 12.313] \\
FD003 & 100 & 10.684 & [7.971, 13.212] \\
FD004 & 248 & 11.192 & [9.253, 13.206] \\
\bottomrule
\end{tabular}
\end{table}

The intervals are wide---4.1 cycles on FD001---because they are honest about
engine-level sampling. Differences between methods smaller than about one cycle on
a 100-engine fleet are not resolvable at this sample size.

\subsection{Uncertainty quantification}

A point estimate is not enough for maintenance scheduling, so we also measure
whether the ensemble knows when it is wrong. Two interval constructions are
compared: raw ensemble sigma, $\hat y\pm z\hat\sigma$ with $\hat\sigma$ the
sample standard deviation over the fifteen models, and split conformal
calibration (CQR) whose half-width is the $(1-\alpha)$ quantile of absolute
residuals on the held-out folds.

Raw ensemble intervals under-cover badly: a nominal 90\% interval covers
61--74\% of engines, and a nominal 95\% interval covers 69--80\%. Fifteen
models that share an architecture, a dataset, and a training protocol make
highly correlated errors, so their spread systematically understates the true
error. CQR restores coverage to 0.82--0.88 at the 90\% level and 0.90--0.94 at
95\%, but the intervals are 25--31 cycles wide at 90\% and 34--45 at 95\%,
against a target range of 0--125. The cause is the heavy tail already visible
in the bias analysis: the largest single error exceeds 60 cycles.

The honest reading is that this model is strong on point accuracy and weak on
uncertainty. A 25-cycle interval does not support a maintenance decision. We
report the numbers rather than a calibration trick that hides them.

\subsection{Instance normalisation as an ensemble-diversity device}

One further ablation is worth reporting because its sign is determined by a
property of the fleet rather than by tuning. Adding a fourth input channel,
the window differenced from its own engine baseline and divided by that
engine's standard deviation, leaves single-model accuracy essentially
unchanged (FD001 $-0.50\%$, FD003 $+0.15\%$, FD002 $+1.46\%$ on official
test) but decorrelates the errors. Five such members added to the fifteen-model
ensemble lower FD001 from 9.64 to 9.47 and FD003 from 10.68 to 10.60.

The same channel is catastrophic on FD004 (validation RMSE 14.85 against 9.5
for the baseline) and does not help FD002. The reason is mechanical: on fleets
with six operating conditions the per-engine variance is dominated by condition
switching, which is an order of magnitude larger than the degradation drift, so
dividing by it compresses precisely the signal the model needs. On
single-condition fleets the per-engine variance is mostly inter-unit nuisance,
and removing it helps. The decision rule is therefore conditioned on the number
of operating conditions, which is known before any test data is seen.

\section{Discussion}

The result we did not expect is that supervising two spans helps while fusing them
hurts. A natural explanation is that the short head acts as a curriculum: it forces
the encoder to keep recent-step information legible, which the long-head loss alone
would allow to decay. At inference the long head already contains that information,
and adding the short head separate estimate only injects its larger variance.

If that reading is right, the design rule generalises: use auxiliary spans as
supervision, never as ensemble members. We have not tested it beyond C-MAPSS, and
the two fleets where the effect is largest are also the ones with the most
operating-condition variation.

\section{Conclusion}

A shared encoder, a distance-from-end positional table, and masked mean pooling
over the full valid history produce official-test RMSE of 9.64, 10.75, 10.68, and
11.19 cycles on C-MAPSS FD001--FD004 with fifteen-model ensembles, 7.6--20.7\% below
a 2026 state-of-the-art baseline on three fleets and 1.9\% above it on the fourth.
The ablation that produced this configuration also produced two negative results:
fusing short and long spans costs 2--6\%, and a redundant second branch costs
18--39\%. The remaining error is concentrated as a positive bias of 5--10 cycles for
engines at 60--90 cycles of remaining life, which is the interval where maintenance
decisions are hardest.

\begin{thebibliography}{9}

\bibitem{saxena2008}
A.~Saxena, K.~Goebel, D.~Simon, and N.~Eklund, Damage propagation modeling for
aircraft engine run-to-failure simulation, in \emph{Proc. Int. Conf. Prognostics}
\emph{and Health Management}, 2008, pp. 1--9.

\bibitem{mfssc2024}
Multi-scale fusion and structural consistency network for aero-engine remaining
useful life prediction, \emph{IEEE Trans. Instrum. Meas.}, 2024.

\bibitem{ast2026}
Decision-conditioned causal graph with state-space backbone for remaining useful
life estimation, 2026.

\bibitem{dkdgjln2026}
Domain-knowledge-guided joint learning network for cross-condition remaining useful
life prediction, 2026.

\end{thebibliography}

\end{document}